\documentclass[11pt]{article}

\usepackage[margin=1in]{geometry}
\usepackage[T1]{fontenc}
\usepackage[utf8]{inputenc}
\usepackage{lmodern}
\usepackage{microtype}
\usepackage{amsmath,amssymb,amsthm,mathtools}
\usepackage{booktabs}
\usepackage{enumitem}
\usepackage{tabularx}
\usepackage{xcolor}
\usepackage[hidelinks]{hyperref}
\usepackage[nameinlink,capitalise]{cleveref}

\newtheorem{theorem}{Theorem}[section]
\newtheorem{proposition}[theorem]{Proposition}
\newtheorem{corollary}[theorem]{Corollary}
\newtheorem{lemma}[theorem]{Lemma}
\theoremstyle{definition}
\newtheorem{definition}[theorem]{Definition}
\newtheorem{assumption}[theorem]{Assumption}
\newtheorem{example}[theorem]{Example}
\theoremstyle{remark}
\newtheorem{remark}[theorem]{Remark}

\newcolumntype{Y}{>{\raggedright\arraybackslash}X}
\newcolumntype{L}[1]{>{\raggedright\arraybackslash}p{#1}}

\newcommand{\cD}{\mathcal D}
\newcommand{\cH}{\mathcal H}
\newcommand{\cK}{\mathcal K}
\newcommand{\cM}{\mathcal M}
\newcommand{\cO}{\mathcal O}
\newcommand{\cR}{\mathcal R}
\newcommand{\cS}{\mathcal S}
\newcommand{\CC}{\mathbb C}
\newcommand{\RR}{\mathbb R}
\newcommand{\1}{\mathbf 1}
\newcommand{\dist}{\operatorname{dist}}
\newcommand{\spec}{\operatorname{spec}}
\newcommand{\Ran}{\operatorname{Ran}}
\newcommand{\Tr}{\operatorname{Tr}}
\newcommand{\dd}{\mathrm d}
\newcommand{\MTT}{Modal Triplet Theory}

\title{\vspace{-1em}
Local Coherent-Sector Robustness Across Model Families:\\[0.25em]
\large Transported Resolvents, Accumulated Error, and the Limits of Universality in MTT}
\author{Peter Nero}
\date{Version 2, July 2026}

\begin{document}
\maketitle

\begin{abstract}
Comparing coherent sectors across different geometries is not the same
problem as proving that one projector is stable inside one Hilbert space.
The models may have different domains, spectra, retained spaces,
observables, clocks, and physical interpretations.  A claim of
``universality'' is therefore meaningful only after those differences are
transported into a common comparison record.

This paper defines that record.  Each model supplies a typed
coherent-sector reduction certificate, a nonempty common spectral region,
and an explicit unitary or quasi-unitary identification of its retained
space with a reference space.  The compressed resolvent is then transported
to the reference space.  We prove a local comparison theorem, a chainwise
error-accumulation bound, and an observable-transfer estimate.  We also
separate projector stability, norm-resolvent robustness, domain or form
convergence, finite-time dynamics, observable agreement, and physical
equivalence.

Two counterexamples mark the boundary of the result.  A continuum of
microscopically distinct operators can have the same retained resolvent and
uniform gap, so local robustness does not make a landscape finite or select
one vacuum.  Conversely, identical reduced spectral data can be paired with
different observable maps and therefore different predictions.  The
current Modal Triplet Theory applications to Calabi--Yau geometry,
heterotic flux backgrounds, perturbative strings, M-theory, and quantum
gravity are consequently recorded as conditional realization or
compatibility programs, not as members of one proved global physical
equivalence class.  The result is a rigorous local comparison framework and
a concrete completion contract for future cross-model claims.
\end{abstract}

\section*{Revision note: Version 2}
\addcontentsline{toc}{section}{Revision note: Version 2}

\begin{description}[leftmargin=2.6cm,style=nextline]
\item[Supersedes]
Version 2 supersedes Version 1.0, DOI
\href{https://doi.org/10.5281/zenodo.18260834}
{10.5281/zenodo.18260834}.

\item[Reason]
Version 1 defined a broad coherent universality class from generic gap,
bounded-geometry, and projector language, then treated Calabi--Yau,
heterotic, string, M-theory, and quantum-gravity constructions as members.
Its projector theorem did not declare the operator family, contour, or
perturbation norm.  Its curvature and RG statements did not define an
evolution or error topology.  It also promoted local spectral control to
full quantum robustness, physical equivalence, and landscape reduction.

\item[Resolution]
This revision imports the domain-explicit Feshbach and projector theorem
from the companion reduction paper rather than duplicating it.  It builds
the missing cross-model layer: explicit comparison spaces, transported
compressed resolvents, accumulated errors, state and observable maps, and
separate dynamical and physical-equivalence gates.

\item[Retained content]
The useful intuition survives locally.  Families of models can share a
controlled effective sector when each model has a valid reduction
certificate and the resulting bounded objects are close after a declared
identification.

\item[Remaining boundary]
No current theorem places all viable compactifications, string limits,
M-theory backgrounds, quantum-gravity constructions, or Standard Model
realizations in one physical equivalence class.  Nor does a local
robustness certificate count, select, or empirically identify vacua.
\end{description}

\tableofcontents

\section{The corrected question}
\label{sec:question}

The original paper asked whether microscopic geometric changes leave one
coherent sector physically unchanged.  That sentence contains three
different mathematical questions:
\begin{enumerate}[leftmargin=2em]
\item Does a selected spectral subspace persist under a perturbation?
\item Can effective operators from different models be compared after
      their retained spaces are identified?
\item Does agreement of those effective operators imply agreement of
      states, observables, dynamics, and physical interpretation?
\end{enumerate}

The first question belongs to spectral perturbation theory.  The second
requires a transported comparison framework, especially when Hilbert
spaces or domains vary.  The third requires additional physical maps.  A
proof of one is not a proof of the next.

\subsection{Why a common gap is not enough}

Suppose two self-adjoint operators each have an isolated low spectral
cluster.  Even then:
\begin{itemize}[leftmargin=2em]
\item their low spaces need not come with a canonical identification;
\item their operators may have different domains;
\item operator-norm, graph-norm, form, strong-resolvent, and semigroup
      convergence are inequivalent hypotheses;
\item their clocks and observable algebras may be different; and
\item the same reduced spectrum can carry different physical labels.
\end{itemize}
The comparison object must therefore include more than a projector and a
gap.

\subsection{What this paper does not claim}

This paper does not prove that:
\begin{itemize}[leftmargin=2em]
\item one coherent projector exists across every \MTT{} encoding;
\item all compactifications lie in one connected or controlled family;
\item local spectral robustness gives uniform long-time dynamics;
\item a finite or Galerkin comparison converges to a continuum theory;
\item reduced operator agreement implies physical equivalence;
\item an admissibility filter makes a candidate landscape finite; or
\item one selected vacuum follows from the existence of a nonempty
      admissible set.
\end{itemize}

\section{The imported single-model reduction}
\label{sec:imported}

The companion paper
\cite{MTTCoherentReduction2026} owns the single-model reduction theorem.
Its essential distinction is worth stating before any cross-model
comparison.

Let
\[
  \cH_a=\cH_{P,a}\oplus\cH_{Q,a}
\]
and let
\[
 T_a=
 \begin{pmatrix}
 A_a&B_a\\ C_a&D_a
 \end{pmatrix},
 \qquad
 \cD(T_a)=\cD(A_a)\oplus\cD(D_a),
\]
where \(A_a\) and \(D_a\) are self-adjoint, \(B_a\) and \(C_a\) are
bounded, and \(C_a=B_a^*\).  For
\(z\in\rho(D_a)\), define
\[
 F_a(z)=A_a-z-B_a(D_a-z)^{-1}C_a.
\]
If \(z\in\rho(T_a)\) as well, then
\[
 R_a^P(z)
 :=
 P_a(T_a-z)^{-1}P_a\big|_{\cH_{P,a}}
 =
 F_a(z)^{-1}.
\]

The bounded correction is
\[
 \Sigma_a(z)=B_a(D_a-z)^{-1}C_a
\]
and obeys
\[
 \|\Sigma_a(z)\|
 \leq
 \|B_a\|\,\|(D_a-z)^{-1}\|\,\|C_a\|.
\]
Every factor is typed and every spectral parameter is declared.

\begin{remark}[Exact spectral projectors]
If \(P_a\) is the exact spectral projector of the same self-adjoint
operator \(T_a\), then it reduces \(T_a\), so
\[
 P_aT_aQ_a=Q_aT_aP_a=0.
\]
A nonzero eliminated-sector correction requires a projector selected from
a reference operator, symmetry, Galerkin space, localization rule, or
approximate criterion.  Cross-model comparison must record which case is
being used.
\end{remark}

This paper begins only after each model has supplied such an individual
certificate.  It does not weaken or silently fill a missing single-model
domain, resolvent, or source row.

\section{The transported comparison record}
\label{sec:record}

\begin{definition}[Controlled model record]
\label{def:modelrecord}
A controlled model record indexed by \(a\) is
\[
\cM_a=
\bigl(
\cH_a,G_a,T_a,P_a,Q_a,
A_a,B_a,C_a,D_a,
\Omega,
\varepsilon_a,
U_a,
\cS_a,\cO_a,\iota_a,
\pi_a
\bigr),
\]
with the following data:
\begin{enumerate}[leftmargin=2em]
\item \(G_a\) is the reference operator or rule selecting \(P_a\);
\item \(T_a\) and its four blocks satisfy the domain assumptions in
      \cref{sec:imported};
\item \(\Omega\) is a nonempty compact set contained in
      \[
      \rho(T_a)\cap\rho(D_a);
      \]
\item the compressed resolvent \(R_a^P(z)\) is uniformly bounded on
      \(\Omega\);
\item \(U_a:\cH_{P,a}\to\cK\) is a declared unitary identification with
      one reference Hilbert space \(\cK\);
\item \(\varepsilon_a\) is the individual reduction or numerical error;
\item \(\cS_a\), \(\cO_a\), and \(\iota_a\) record states, observables,
      and physical interpretation maps; and
\item \(\pi_a\) records provenance, theorem status, and source hashes.
\end{enumerate}
\end{definition}

The common \(\Omega\) is part of the comparison.  Two reductions evaluated
at unrelated spectral parameters or on disjoint resolvent regions have not
yet been compared.

\begin{definition}[Transported compressed resolvent]
For a controlled model record, define
\[
 \widehat R_a(z)
 =
 U_a R_a^P(z)U_a^*
 \in\mathcal B(\cK),
 \qquad z\in\Omega.
\]
For two records define
\[
 d_\Omega(a,b)
 =
 \sup_{z\in\Omega}
 \|\widehat R_a(z)-\widehat R_b(z)\|.
\]
\end{definition}

Because the identifications \(U_a\) are part of the records, this
functional compares bounded operators on one fixed space.  Changing
\(U_a\) changes the question and must be reported as a different
comparison convention.

\section{Local comparison and accumulated error}
\label{sec:comparison}

\begin{theorem}[Transported local comparison]
\label{thm:transported}
Let \(a,b\) be controlled model records on a common compact region
\(\Omega\).  Then \(d_\Omega(a,b)<\infty\).  It is symmetric and obeys
the triangle inequality.  More generally, suppose
\[
 E_a,E_b:\Omega\to\mathcal B(\cK)
\]
are declared approximate reduced-resolvent fields satisfying
\[
 \sup_{z\in\Omega}\|\widehat R_a(z)-E_a(z)\|
 \leq\varepsilon_a,
 \qquad
 \sup_{z\in\Omega}\|\widehat R_b(z)-E_b(z)\|
 \leq\varepsilon_b.
\]
Then
\[
 d_\Omega(a,b)
 \leq
 \varepsilon_a
 +
 \sup_{z\in\Omega}\|E_a(z)-E_b(z)\|
 +
 \varepsilon_b.
\]
\end{theorem}

\begin{proof}
Uniform boundedness follows from compactness of \(\Omega\), continuity of
each resolvent, and the unitary invariance of the operator norm.  Symmetry
is immediate.  For any third record \(c\),
\[
\begin{aligned}
\|\widehat R_a(z)-\widehat R_c(z)\|
&\leq
\|\widehat R_a(z)-\widehat R_b(z)\|\\
&\quad+
\|\widehat R_b(z)-\widehat R_c(z)\|.
\end{aligned}
\]
Taking the supremum gives the triangle inequality.  Inserting and
subtracting \(E_a(z)\) and \(E_b(z)\) gives the final estimate.
\end{proof}

\begin{corollary}[Error accumulation along a comparison chain]
\label{cor:chain}
For controlled records \(a_0,\ldots,a_N\) on the same \(\Omega\),
\[
 d_\Omega(a_0,a_N)
 \leq
 \sum_{j=0}^{N-1}d_\Omega(a_j,a_{j+1}).
\]
Consequently, a chain of \(N\) comparisons each certified only to
\(\varepsilon\) gives at best an \(N\varepsilon\) endpoint bound unless a
stronger cancellation theorem is supplied.
\end{corollary}

\begin{proof}
Iterate the triangle inequality in
\cref{thm:transported}.
\end{proof}

This simple bound matters in a large corpus.  A sequence
\[
 \text{upper geometry}
 \longrightarrow
 \text{finite operator}
 \longrightarrow
 \text{four-dimensional action}
 \longrightarrow
 \text{observable}
\]
cannot reuse the same tolerance at every arrow and report it once at the
end.  Each arrow contributes its own error and provenance row.

\subsection{Local class, not global equivalence}

\begin{definition}[Local coherent comparison class]
Fix a reference record \(a_0\), a common region \(\Omega\), fixed
identification conventions, and a tolerance \(\varepsilon>0\).  The local
comparison class is
\[
 \mathfrak U_\varepsilon(a_0;\Omega)
 =
 \{a:\ d_\Omega(a,a_0)\leq\varepsilon\}.
\]
\end{definition}

This is a useful neighborhood, not a theorem that all viable models lie in
it.  Finite-tolerance membership is also not automatically transitive at
the same tolerance: two members can be as far as \(2\varepsilon\) apart.

\section{Projector stability is one separate layer}
\label{sec:projector}

The retained spaces themselves may vary.  To compare their projectors, one
needs full-space identifications, not only the retained unitaries in
\cref{def:modelrecord}.

\begin{proposition}[Transported Riesz-projector bound]
\label{prop:riesz}
Let \(W_a:\cH_a\to\cH_*\) be unitary and set
\[
 \widetilde G_a=W_aG_aW_a^*.
\]
Suppose \(G_*\) and \(\widetilde G_a\) are self-adjoint on the same dense
domain,
\[
 V_a=\widetilde G_a-G_*
\]
is bounded and self-adjoint, and a positively oriented contour
\(\Gamma\subset\rho(G_*)\) encloses the selected reference cluster.  Define
\[
 M_*=\sup_{z\in\Gamma}\|(G_*-z)^{-1}\|.
\]
If \(M_*\|V_a\|<1\), then \(\Gamma\subset\rho(\widetilde G_a)\), and the
corresponding spectral projectors satisfy
\[
 \|W_aP_aW_a^*-P_*\|
 \leq
 \frac{\operatorname{length}(\Gamma)}{2\pi}
 \frac{M_*^2\|V_a\|}
      {1-M_*\|V_a\|}.
\]
If the right-hand side is below one and \(P_*\) has finite rank, the
retained ranks agree.
\end{proposition}

\begin{proof}
This is the standard Riesz-projector perturbation estimate in the
transported Hilbert space.  The resolvent identity and Neumann series give
\[
 \|(\widetilde G_a-z)^{-1}\|
 \leq
 \frac{M_*}{1-M_*\|V_a\|}.
\]
Integrating the second resolvent identity around \(\Gamma\) yields the
bound.  Orthogonal projectors at distance below one have isomorphic
ranges.  See Kato and Davis--Kahan
\cite{Kato1995,DavisKahan1970}; the domain-explicit \MTT{} specialization
is recorded in \cite{MTTCoherentReduction2026}.
\end{proof}

\begin{remark}
\Cref{prop:riesz} proves projector and finite-rank stability.  It does not
prove closeness of the full operators, their Feshbach maps, semigroups,
states, observables, or physical interpretations.
\end{remark}

\section{When the Hilbert spaces genuinely vary}
\label{sec:varying}

A unitary \(W_a:\cH_a\to\cH_*\) always exists for Hilbert spaces of the same
Hilbert dimension, but it may not represent the geometry or preserve the
operators.  Singular limits can also change effective spaces.  The
appropriate literature then uses identification maps and explicit defect
bounds rather than pretending the operators already act on one space
\cite{KuwaeShioya2003,Post2012}.

\begin{definition}[Quasi-unitary comparison data]
A quasi-unitary comparison from \(\cH_{P,a}\) to \(\cK\) consists of
\[
 J_a:\cH_{P,a}\to\cK,
 \qquad
 J_a':\cK\to\cH_{P,a},
\]
and a defect \(\eta_a\geq0\) such that
\[
 \|J_a\|,\|J_a'\|\leq1+\eta_a,
\]
\[
 \|J_a'J_a-\1_{\cH_{P,a}}\|\leq\eta_a,
 \qquad
 \|J_aJ_a'-\1_{\cK}\|\leq\eta_a.
\]
One then transports the compressed resolvent as
\[
 \widetilde R_a(z)=J_aR_a^P(z)J_a'.
\]
\end{definition}

The defect must enter the error budget.  For example, if
\(\|R_a^P(z)\|\leq r_a\) on \(\Omega\), then
\[
 \|\widetilde R_a(z)\|
 \leq
 (1+\eta_a)^2r_a.
\]
More substantive convergence statements require compatibility of the
forms, domains, or resolvents with \(J_a,J_a'\).  Mosco-type convergence
and varying-space norm-resolvent convergence are established frameworks
for such problems \cite{KuwaeShioya2003,Post2012}.  The phrase ``same
coherent sector'' is not a replacement for those hypotheses.

\section{Seven robustness layers}
\label{sec:layers}

\begin{table}[htbp]
\centering
\small
\begin{tabularx}{\textwidth}{L{0.20\textwidth}L{0.28\textwidth}Y}
\toprule
Layer & Typical certificate & What it does not imply \\
\midrule
Selected rank
& Riesz contour, perturbation norm, rank bound
& Close reduced operator \\
Projector
& Operator-norm or angle estimate after full-space transport
& Stable domains or dynamics \\
Reduced resolvent
& Common \(\Omega\), retained identification, norm bound
& Uniform long-time evolution \\
Domain or form
& Graph-gap, relative-bound, or Mosco hypotheses
& Observable agreement \\
Finite-time dynamics
& Generator or semigroup estimate on a declared interval
& Arbitrary-time equivalence \\
States and observables
& State channel, observable transport, expectation error
& Same microscopic theory \\
Physical equivalence
& Algebra, states, dynamics, symmetries, clocks, and comparison functor
& Follows from no single spectral row \\
\bottomrule
\end{tabularx}
\caption{Robustness is a stack of distinct claims, not one adjective.}
\label{tab:layers}
\end{table}

\subsection{Finite-time dynamics}

\begin{proposition}[Bounded-generator finite-time comparison]
\label{prop:duhamel}
Let \(H_a,H_b\) be bounded self-adjoint operators on \(\cK\) and suppose
\[
 \|H_a-H_b\|\leq\delta.
\]
Then for every \(t\in\RR\),
\[
 \|e^{-itH_a}-e^{-itH_b}\|
 \leq |t|\delta.
\]
\end{proposition}

\begin{proof}
Duhamel's formula gives
\[
 e^{-itH_a}-e^{-itH_b}
 =
 -i\int_0^t
 e^{-i(t-s)H_a}(H_a-H_b)e^{-isH_b}\,\dd s.
\]
The unitary factors have norm one.
\end{proof}

This is a standard functional-analytic estimate, not a new \MTT{} law.  It
also displays the time horizon explicitly.  A local spectral or generator
error can accumulate into an order-one phase error at sufficiently long
times.

\subsection{States and observables}

\begin{proposition}[Expectation-transfer bound]
\label{prop:observable}
Let \(\rho_a,\rho_b\) be density operators on \(\cK\), and let
\(O_a,O_b\in\mathcal B(\cK)\) be bounded observables.  Then
\[
\left|
\Tr(\rho_aO_a)-\Tr(\rho_bO_b)
\right|
\leq
\|\rho_a-\rho_b\|_1\,\|O_a\|
+
\|O_a-O_b\|.
\]
\end{proposition}

\begin{proof}
Insert and subtract \(\Tr(\rho_bO_a)\).  Trace duality bounds the first
difference by
\(\|\rho_a-\rho_b\|_1\|O_a\|\).  Since
\(\|\rho_b\|_1=1\), the second is bounded by
\(\|O_a-O_b\|\).
\end{proof}

The proposition shows why reduced-resolvent agreement is not yet empirical
agreement.  One must also transport states and observables and bound their
errors.

\section{Three boundary counterexamples}
\label{sec:counterexamples}

\begin{example}[A robust continuum need not be a small landscape]
\label{ex:continuum}
For each \(\lambda\in[0,1]\), let
\[
 \cH_\lambda=\CC^3,
 \qquad
 T_\lambda=
 \begin{pmatrix}
 0&0&0\\
 0&2&0\\
 0&0&3+\lambda
 \end{pmatrix},
 \qquad
 P_\lambda=
 \begin{pmatrix}
 1&0&0\\
 0&0&0\\
 0&0&0
 \end{pmatrix}.
\]
The selected eigenvalue is isolated by a uniform gap of at least two.
All retained spaces are the same line, and for every common
\(z\neq0\),
\[
 R_\lambda^P(z)=-z^{-1}.
\]
Thus
\[
 d_\Omega(\lambda,\mu)=0
\]
for all \(\lambda,\mu\), while the eliminated spectrum varies continuously.
Local coherent-sector robustness is therefore compatible with a continuum
of microscopically distinct models.
\end{example}

\begin{example}[The same reduced spectrum can give different observables]
\label{ex:observable}
Take two records with the same one-dimensional retained Hilbert space and
the same compressed resolvent.  Give the first record the observable
\(O_0=0\) and the second \(O_1=\1\).  Their unique normalized state gives
\[
 \Tr(\rho O_0)=0,
 \qquad
 \Tr(\rho O_1)=1.
\]
The spectral comparison error is zero, but the predictions differ.  An
observable identification is indispensable.
\end{example}

\begin{example}[Projector stability is not arbitrary-time stability]
\label{ex:time}
On \(\CC\), take \(H_0=0\) and \(H_\varepsilon=\varepsilon\).  The projector
is identical and the resolvents converge away from their poles.  At
\[
 t=\frac{\pi}{\varepsilon},
\]
however,
\[
 \|e^{-itH_\varepsilon}-e^{-itH_0}\|=2.
\]
Finite-time bounds do not become uniform-in-time equivalence merely because
\(\varepsilon\) is small.
\end{example}

\section{What landscape reduction would require}
\label{sec:landscape}

Let \(\mathfrak V\) be a declared candidate set and let
\[
 C:\mathfrak V\to\{0,1\}
\]
be a certified admissibility predicate.  The surviving set is
\[
 \mathfrak V_{\rm adm}
 =
 \{v\in\mathfrak V:C(v)=1\}.
\]
This is a filter.  It becomes a landscape-reduction result only after
additional questions are answered:
\begin{enumerate}[leftmargin=2em]
\item How is \(\mathfrak V\) defined or enumerated?
\item Is \(C\) decidable or certifiable on every candidate in scope?
\item What measure, index, or counting convention defines ``most''?
\item Is the surviving set finite, discrete, compact, or merely nonempty?
\item What dynamics or selection rule chooses among multiple survivors?
\item Which survivor is empirically identified with our universe?
\end{enumerate}

Flux-vacuum counts in string theory illustrate why these inputs matter:
they start from a specified ensemble, charge cutoff, and index or
distribution \cite{AshokDouglas2004,DouglasKachru2007}.  A local operator
gap by itself provides none of that counting data.

\begin{proposition}[An admissibility predicate does not imply uniqueness]
\label{prop:nonselection}
If \(|\mathfrak V_{\rm adm}|>1\), then the predicate \(C\) alone does not
select a unique element of \(\mathfrak V_{\rm adm}\).
\end{proposition}

\begin{proof}
Every survivor has the same predicate value \(1\).  No rule using only that
value distinguishes two different survivors.
\end{proof}

The logical separation among mathematical existence, chart validity,
admissibility, realization, and empirical identification is developed in
\cite{MTTVacuumSelection2026}.  Local robustness belongs between
admissibility and controlled realization; it is not a substitute for the
later layers.

\section{Current MTT application map}
\label{sec:applications}

The present paper does not use older papers as proof sources merely because
they contain the word ``universality.''  It reads the selected current
revisions at their declared tiers.

\begin{table}[htbp]
\centering
\small
\begin{tabularx}{\textwidth}{L{0.20\textwidth}L{0.25\textwidth}Y}
\toprule
Target family & Current established role & Missing comparison or selection rows \\
\midrule
Fixed-point framework
& Conditional analytic existence and stability under stated hypotheses
& Physical carrier, source map, and realization theorem
  \cite{MTTFixedPointsI2026} \\
Calabi--Yau
& Conditional lower realization target; fixed compact elliptic problems have standard spectral control
& Uniformity over moduli, stable visible--hidden bundle, worldsheet, stabilization, and unique selection
  \cite{MTTCalabiYau2026} \\
Heterotic/Fu--Yau
& Exact local and finite ingredients; Lens--Nil is auxiliary; selected \(q=79\) is a completion program
& Physical rank-three visible--hidden endpoint and common comparison operator
  \cite{MTTHeteroticFlux2026,MTTHeteroticSelection2026} \\
Perturbative string
& Conditional typed encoding at declared order and genus
& Complete physical worldsheet record, exact infrared SCFT, GSO/modular data, and quantum completion
  \cite{MTTPerturbativeString2026} \\
M-theory low-energy limit
& Conditional eleven-dimensional target record
& Selected spin geometry, differential C-field, action, branes, and action-preserving upper map
  \cite{MTTMTheory2026} \\
Quantum gravity
& Exact finite TT support and conditional/fixed-order low-energy parity
& Physical normalization, complete state and measure, selected higher-derivative values, and nonperturbative UV completion
  \cite{MTTQuantumGravity2026} \\
\bottomrule
\end{tabularx}
\caption{Current branches are not one proved global physical equivalence class.}
\label{tab:applications}
\end{table}

This table does not say the programs are incompatible.  It says that a
common comparison class must be earned by constructing the same typed rows
for each member and transporting them through explicit maps.

\section{A complete cross-model certificate}
\label{sec:certificate}

\begin{definition}[Cross-model robustness certificate]
A cross-model certificate for records \(a,b\) contains
\[
\mathfrak C_{a\leftrightarrow b}
=
\bigl(
\mathfrak C_a,\mathfrak C_b,
\Omega,
U_a,U_b,
d_\Omega,
\varepsilon_{\rm chain},
\mathfrak T,
\mathfrak O,
\mathfrak I,
\pi
\bigr),
\]
where:
\begin{itemize}[leftmargin=2em]
\item \(\mathfrak C_a,\mathfrak C_b\) are complete single-model reduction
      certificates;
\item \(\Omega\) is the common full and eliminated-sector resolvent region;
\item \(U_a,U_b\), or their quasi-unitary replacements, are the
      identification maps;
\item \(d_\Omega\) is the transported compressed-resolvent error;
\item \(\varepsilon_{\rm chain}\) is the accumulated error across all
      intermediate arrows;
\item \(\mathfrak T\) records any finite-time, semigroup, graph, or form
      comparison;
\item \(\mathfrak O\) records state, observable, and expectation transport;
\item \(\mathfrak I\) records clocks, units, symmetries, and physical
      interpretation; and
\item \(\pi\) records provenance and theorem status.
\end{itemize}
\end{definition}

\subsection{Numerical execution}

For a finite or Galerkin comparison, an executable packet should report:
\begin{enumerate}[leftmargin=2em]
\item basis dimensions and immutable basis hashes;
\item Hermiticity and projector residuals;
\item source and rank of every projector;
\item interval enclosures for the common spectral region;
\item certified norms of the eliminated resolvents and couplings;
\item the retained identification matrices and their unitary defects;
\item interval bounds for \(d_\Omega\);
\item every discretization, reduction, transport, and roundoff error
      separately;
\item state and observable comparison rows when empirical claims are made;
      and
\item a continuum-transfer theorem when the physical target is not the
      finite object itself.
\end{enumerate}

An exact finite calculation is exact for the finite object.  It is not a
continuum, string, gravitational, or physical-equivalence theorem until
the corresponding transfer rows are supplied.

\section{Ownership of the mathematics}
\label{sec:ownership}

Spectral projections, resolvent perturbation, closed-operator stability,
and semigroup estimates belong to established operator theory
\cite{Kato1995,DavisKahan1970}.  Frameworks for convergence across varying
Hilbert spaces and spectral structures are likewise established
\cite{KuwaeShioya2003,Post2012}.

The \MTT{} Foundation owns the general admissibility and status semantics
\cite{MTTFoundation2026}.  The companion coherent-sector paper owns the
domain-explicit Feshbach reduction, exact-projector correction, and
single-model robustness estimate \cite{MTTCoherentReduction2026}.  The
current Calabi--Yau, heterotic, string, M-theory, and quantum-gravity papers
own their individual realization records and boundaries.

This paper owns only the integration layer:
\begin{itemize}[leftmargin=2em]
\item the transported cross-model record;
\item the common-region compressed-resolvent comparison;
\item chainwise accumulation of reduction and comparison errors;
\item the state-observable transfer requirement;
\item the explicit seven-layer robustness hierarchy; and
\item the counterexamples separating local robustness from landscape
      reduction and physical equivalence.
\end{itemize}

\section{Claim status and limitations}
\label{sec:status}

\begin{table}[htbp]
\centering
\small
\begin{tabularx}{\textwidth}{L{0.31\textwidth}L{0.18\textwidth}Y}
\toprule
Statement & Status here & Boundary \\
\midrule
Single-model Feshbach reduction
& Imported
& Domain-explicit companion theorem \\
Transported resolvent comparison
& Proved
& Common \(\Omega\), declared identifications \\
Chainwise error accumulation
& Proved
& Sum bound; no assumed cancellation \\
Riesz-projector stability
& Imported and specialized
& Fixed transported domain, bounded perturbation \\
Finite-time bounded dynamics
& Standard and proved
& Bounded generators, declared time interval \\
Expectation transfer
& Proved
& Density operators and bounded observables \\
Varying-space convergence
& Contract only
& Requires quasi-unitary, Mosco, form, or resolvent theorem \\
One global MTT universality class
& Not established
& No common source and full comparison record \\
Landscape reduction
& Not established
& Requires candidate set, count/measure, filter execution, and selection \\
Physical equivalence
& Not established
& Requires states, observables, dynamics, symmetries, clocks, and units \\
\bottomrule
\end{tabularx}
\caption{Claim status in Version 2.}
\label{tab:status}
\end{table}

\section{Completion and falsifiability contract}
\label{sec:completion}

A claimed cross-model class is complete only if it supplies:
\begin{enumerate}[leftmargin=2em]
\item one exact list of models or a precisely quantified family;
\item one valid single-model reduction certificate for every member;
\item explicit full-space or retained-space identification maps;
\item a nonempty common spectral region;
\item fixed-domain, graph, form, or varying-space hypotheses appropriate
      to the claimed topology;
\item a certified transported error and an accumulated chain budget;
\item finite-time or semigroup control if dynamics is claimed;
\item state and observable maps if predictions are compared;
\item clock, unit, symmetry, and interpretation maps if physical
      equivalence is claimed; and
\item an enumeration, measure, decision procedure, and selector if
      landscape reduction or unique-vacuum selection is claimed.
\end{enumerate}

The claim fails in its declared regime if the contour meets the spectrum,
the retained rank changes, the common resolvent region becomes empty, an
identification defect exceeds its tolerance, accumulated error exceeds the
observable budget, the domain topology changes without control, or a
held-out observable violates the expectation bound.

\section{Conclusion}
\label{sec:conclusion}

Coherent-sector robustness is a meaningful and useful local statement.
It says that after each model has supplied a valid reduction and after the
retained spaces have been explicitly identified, their bounded compressed
resolvents can be compared with a certified error.

That statement does not become global universality by changing its name.
Projectors, domains, resolvents, finite-time dynamics, states, observables,
and physical interpretation are different layers.  Their errors accumulate
through a chain.  A continuum of microscopic models can share exactly the
same retained resolvent, and equal spectral data can still support
different observables.

The corrected role of this paper is therefore precise.  It supplies the
mathematical language by which future \MTT{} branches can prove that they
belong to one controlled local comparison class.  The current
Calabi--Yau, heterotic, string, M-theory, and quantum-gravity branches are
valuable conditional targets and partial constructions, but their common
physical equivalence and any resulting landscape reduction remain theorem
obligations rather than established conclusions.

\section*{Reproducibility statement}

No numerical result is claimed in this paper.  Every displayed estimate is
analytic and can be checked directly from the stated assumptions.
Application-level computations must publish the single-model certificates,
identification maps, interval or exact norm bounds, source hashes, and
accumulated error ledger described in \cref{sec:certificate}.

\bibliographystyle{plain}
\bibliography{main}

\end{document}
